% Generated draft: requires mathematical and source review.
\documentclass[11pt]{article}
\usepackage[margin=25mm]{geometry}
\usepackage{fontspec}
\setmainfont{Latin Modern Roman}
\usepackage{amsmath,amssymb,mathtools,unicode-math}
\setmathfont{Latin Modern Math}
\usepackage{enumitem,needspace,xurl,hyperref}
\setlength{\parindent}{0pt}
\setlength{\parskip}{5pt}
\setlength{\emergencystretch}{3em}
\newcommand{\N}{\mathbb{N}}
\newcommand{\Z}{\mathbb{Z}}
\newcommand{\Q}{\mathbb{Q}}
\newcommand{\R}{\mathbb{R}}
\newcommand{\C}{\mathbb{C}}
\newcommand{\F}{\mathbb{F}}
\newcommand{\A}{\mathbb{A}}
\newcommand{\PP}{\mathbb{P}}
\newcommand{\E}{\mathbb{E}}
\newcommand{\eps}{\varepsilon}
\newcommand{\dd}{\,\mathrm{d}}
\begin{document}
\section*{UnsolvedMath ID: 30001892}
\textbf{Source problem code:} catalogue entry
\subsection{Definitions and mathematical statement}
A provisional conventional framework is the following. Let \(S\subset\R^3\) be a finite set and let \(r:S\to(0,\infty)\) assign a favourite distance to each point. The associated directed graph has vertex set \(S\) and arc set
\[
A(S,r)=\{(x,y)\in S\times S:x\ne y,\ \|x-y\|_2=r(x)\}.
\]
For a positive integer \(n\), one possible meaning of extremality is that \(|S|=n\) and \(|A(S,r)|\) is maximal over all such pairs with \(n\) points. These definitions provide a provisional interpretation; the supplied record does not specify the admissible configurations or the extremal objective explicitly.

The supplied question asks whether the exceptional set \(T\) in the extremal favourite-distance digraph problem in \(\R^3\) is always empty. Its dated literature assessment describes \(T\) as having at most two points and says that its emptiness is conjectured.

\textbf{Unresolved formulation.} The record does not define \(T\), state the structural result that introduces it, or specify the range of \(n\). It also does not say whether emptiness is required for every permitted structural decomposition or whether every extremal configuration merely admits a decomposition with no exceptional points. Consequently, the intended assertion \(T=\varnothing\) cannot be given complete hypotheses and quantifiers from the supplied evidence. In particular, no theorem asserting a decomposition or identifying its nonexceptional part is assumed here.
\subsection{Short English statement}
The question asks whether extremal favourite-distance digraphs in \(\R^3\) always have no exceptional points. The supplied assessment mentions at most two possible exceptional points, but their definition, the relevant structural decomposition, and the applicable configuration sizes are missing. A precise conjecture therefore requires the underlying source formulation.
\subsection{Sources}
\begin{itemize}
\item Supplied local UnsolvedMath registry record, permanent ID \(30001892\), title \emph{Exceptional Points in Favourite Distance Digraphs}, exact source problem code \texttt{catalogue entry}; collection: Oberwolfach Reports Open Problems. The record lists proposed year 2011 and names no proposer.
\item Supplied report citation: \emph{Discrete Geometry (2012)}, Oberwolfach Reports, DOI \texttt{10.4171/owr/2011/44}, pp.\ 2459--2548. \url{https://doi.org/10.4171/owr/2011/44}. The background describes these reports as documenting mathematical workshops and their open-problem discussions. The literature-source label is ``OWR 44 (2012).''
\item Supplied literature source labelled ``Swanepoel'': \url{https://researchonline.lse.ac.uk/id/eprint/104165/4/Swanepoel_favourite_distances_in_3_space_published.pdf}. The supplied assessment refers to a 2020 primary paper stating that the exceptional set of at most two points is conjectured to be empty; the record supplies no theorem number or relevant passage.
\item Literature checked at: \texttt{2026-08-21}. This is supplied dated curation, not independently verified. The supplied status is ``open,'' and the background explicitly requests independent verification of time-sensitive status and cited claims.
\item The supplied statement status is \texttt{corrected\_verified}; its verification note says only that TeX notation was enclosed in mathematical delimiters without changing the mathematical content. This does not establish completeness or mathematical correctness of the formulation.
\end{itemize}
\end{document}
